\documentclass{jsarticle}
\usepackage{amsmath}
\usepackage{amssymb}
\begin{document}
\title{嶋本淑子さんとの関係から \\
微分幾何の量子化の \\
計算方法がわかり、それが大域的微分多様体を \\
単体的積分と単体的微分へと \\
進化させることにしたきっかけのレポート}
\author{Masaaki Yamaguchi}
\date{}
\maketitle

   ベータ関数をガンマ関数へと渡すと、
   $$ \beta(p,q) = {{\Gamma(p)\Gamma(q)} \over {\Gamma(p + q)}} $$
   ここで、単体的微分と単体的積分を定義すると
   $$ t = \Gamma(x) $$
   $$ T = \int \Gamma(x) dx $$
   これを使うのに、ベータ関数を単体的微分へと渡すために、
   $$ \beta(p,q) = - \int {1 \over {t^2}}dt $$
   大域的微分変数へとするが、
   $$ T_m = \int \Gamma(x) dx_m $$
   この単体的積分を常微分へとやり直すと、
   $$ T^{'} = {t^{'} \over {{\log t}}}dt + C(\text{Cは積分定数}) $$
   これが、大域的微分多様体の微分変数と証明するためには、
   $$ dx_m = {1 \over {\log x}}dx, dx_m = {(\log x^{-1})}^{'}  $$
   これより、単体的微分が大域的部分積分へと置換できて、
   $$ T^{'} = \int \Gamma(\gamma)^{'} dx_m $$
   微分幾何へと大域的積分と大域的微分多様体が、加群分解の共変微分へと
   書き換えられる。
   $$ \bigoplus T^{\nabla}  = \int T dx_m, \delta(t) = t dx_m  $$


   これらをまとめると、ベータ関数を単体的関数として、
   これを常微分へと解を導くと、単体的微分と単体的積分
   へと定義できて、これより、大域的微分多様体と大域的積分多様体
   へとこの関数が存在していることを証明できるのを上の式たちから
   わかる。この単体的微分と単体的積分から微分幾何が書けることがわかる。

   
   これらは自分が導いた方程式たちとおもうようにする。

   
   微分の本質は、差と商の計算であり、積分の本質は加群である。
   この微分と積分を合わせた計算が加群分解であり、ガロワ群である。

   アカシックレコードにいると思っていた私の子どもとも勘違いしていた、
   子どもはまぼろしであり、架空のストーリーの作為体験の子ともと思う次第
   であり、架空のストーリーに載せられているとの妄想癖での、
   このストーリーに載せられている自体の妄想が作為体験である。
   作為があったとは違う。日常の作為と妄想の作為は違う。

   レドックス電池を言っていたから、荒木先生の経営している病院の
   医者である江松先生が私を助けようとして、
   作為体験のために、私のあのときまでの情報、
   私の周りの情報を使って、架空のストーリーを立てて、作為体験として、
   あのときは、ＭＲＩ室の壁を取っていて、壁が無かった。
   先生は、東北の地震があるのを事前に知っていたとしか思えないことを、
   福山市の私の散歩コースに現れて、挨拶とを交わしていて、
   本当に庇ってくれて申し訳ないです。ＳＲＳ速読法をしていたのを、
   占さんがアルバイトのあとに、トポロジーのレポートを見せていたら、
   気づいたらしく、東北のことを知っていたらしく、史記と三国志を
   私に渡した上で、読書力を中心として、勉強をしなさいとフォローを
   してくれて、していたら、２月２３日に持田さんと伊藤一朗さんが
   ラジオ放送に出ていて、シングルをリリースしましたと言うので、
   YouTubeを見ると、一箇所だけ動画がとんでいて、かばうのを
   彩さんにすればいいのに、統合失調症の真逆のうれしいのに悲しむとは
   感情と現状での置かれている立場から違うが、似たような作用の症状でもあり、
   「存在と無」を啓文社に予約したのが、取りに行くのに１８：００
   にしたのが、私の顛末であり、全部私の責任です。７ヶ月間びっしりと
   いたぶられました。増援は蔵王病院らしく、とっくに気づいています。
   史記は、小方君と占さんのおかげで助かって、三国志が私に増援がくると
   助言していた。占さんが蔵王病院に行かしたのは、１０月１日に彩さん
   がめざましテレビを辞めたからと、お父さんの体調を気にしてのことと
   ３年前くらいからわかっていました。２０１３年に彩さんがめざましテレビ
   を辞めたのを知ったが、２０１８年の６月に何時辞めたかを知り、
   携帯サイトでのKindleから１８７７回目のめざましテレビの放送で
   辞めたのを知り、姉が怒っていたのを、いろんないきさつがあったのを
   知り、私の現状を自身にふりかかっている出来事に怒ることが、
   少なくなりました。
   グロス・グラスマンは、間所くん、青木くん、
   伊澤くん、圭一さん、寺林さん、鷺坂くん、下神くん、梶岡くん、
   片寄くん、鶴元くん、圭一くんです。
   鶴元くんの元は、歌舞伎や能を受け継ぐ家系の元であり、家元５代流
   という元である。
   占さんを入れたいが、周りの批判を受けるので、入れることができない。
   渡辺由美さんは、私の恋人として、
   彩さんに継ぐ大切な人でもある。相手が私の身代わりとして結婚したと、
   デイケアのキッチンで言っていた。気づいていたが、私本人と結婚しろよと
   ツッコミを入れたかった。彩さんにも言いたいです。
   大域的微分多様体の基軸になる、基本群をもとめることができて、
   $$ \pi(\chi,x) = [i\pi(\chi,x), f(x)] $$
   
   一般相対性理論の式は、
   $$ \kappa T^{\mu\nu} = R_{\mu\nu} + {1 \over 2}g_{ij}\Lambda $$
   これを大域的微分多様体と同型といえるのは、
   $$  {d \over df}F = {d \over df}\int\int{1 \over {(x \log x)^2}}dx_m
   + {d \over df}\int\int{1 \over {(y \log y)^{1 \over 2}}}dy_m $$
   ここで、一般相対性理論の多様体積分は、
   $$ \int\left( {R_{\mu\nu} + {{1 \over 2}g_{ij}\Lambda}}\right)\mathrm{
   dvol} $$
   この式は、一般相対性理論は、ガウスの曲平面による重力場方程式
   であり、この多様体積分は、世界面である偏微分方程式を大域的に積分して、
   その積分多様体をヘルマンダー作用素と同型に微分すると、
   $$ ||ds^2|| = \nabla \nabla \int \nabla f(x)d \eta $$
   $$  = \int {\kappa T^{\mu\nu}}\mathrm{dvol} = - 2 R_{ij} $$
   これは、大域的微分方程式の重力と反重力方程式と同型であり、
   リッチ・フロー方程式になり、
   $$ = \int\left(R_{\mu\nu}\right)\mathrm{dvol} + \int{{1 \over 2}g_{ij}
   \Lambda}\mathrm{dvol} = \text{重力場方程式 + 反重力場方程式(斥力)} $$
   $$ = {d \over df}\int\int{1 \over {(x \log x)^2}}dx_m
   + {d \over df}\int\int{1 \over {(y \log y)^{1 \over 2}}}dy_m $$
   と、多様体積分が大域的微分多様体となる。
   要するに、世界面での一般相対性理論を多様体積分すると、宇宙において、
   伝播する単位空間における立体面の重力波になるのが、マクロレベルで、
   宇宙のまわりのブラックホールのシュバルツシルト半径であり、
   Jones多項式にもとまる。まとめると、一般相対性理論を多様体積分すると、
   立体面の重力波が宇宙規模の宇宙のまわりのブラックホール解になるということと、
   Jones多項式が統一場理論となることである。


   そうゆう経緯であり、なぜアカシックレコードの私の子どもが
   ２０２３年の１０月に狙いをさだめているかがわかるのは、
   私がご先祖様にアクセスするのと原理が同じだが、ドアーを使う
   アカシックレコードへのアクセスによっても、微分幾何の量子から
   私の大域的微分多様体の計算方法が間接的にわかったことでもあり、
   直接でもある、思念同士の交流でもある方法だからです。
   彩さんが今産みたいとおもっても、あの子は待つと思うという次第です。
   １２月２５日に狙いを定めていて、うちのユウと同じことをしようとしている
   という次第です。彩さんは、マスターベーションで産みたいらしいので、
   私は文句を言わんが、あまり、私をいじめないでほしいです。
   あの子は了解済みらしいです。

  \end{document}
